\chapter{教学矛盾决策表}

本表借用 TRIZ 的“先说清矛盾，再寻找分离或转化条件”思路。它是备课启发工具，不是经过验证的教学理论。

\section*{开放探索 / 专业正确}
常见失败是完全放任或直接给答案。可以先让学生自由生成方案，再设置强制核验关。检查：学生何时必须回到规范、量纲或边界条件？

\section*{即时反馈 / 独立思考}
常见失败是 AI 过早给出答案。可以先提交预测，再开放提示与计算。检查：学生是否在看答案前作过一次可见尝试？

\section*{个性化 / 教师负担}
常见失败是无限制逐份精批。可以让工具检查低层错误，教师聚焦关键判断。检查：教师反馈是否用在只有人能承担的地方？

\section*{游戏吸引 / 学习目标}
常见失败是积分热闹、知识游离。可以让核心操作本身产生学科反馈。检查：得高分是否真的需要使用目标知识？

\section*{作品精美 / 过程真实}
常见失败是最终图像遮蔽错误。可以强制提交初稿、版本差异和口头答辩。检查：能否看见作品如何变成现在的样子？

使用时只选一对矛盾。先写清当前课程在哪一端失衡，再设计一个最小改动，并用学生作品验证，而不是一次引入所有新工具。
